\documentclass[12pt]{article}
%% BUGS:
%% - The \begin{problem} ... \end{problem} is only used in exams??

\usepackage{url, graphicx, epstopdf}

% page layout
\usepackage[letterpaper, textheight=9.5in, textwidth=5.5in]{geometry}
\usepackage{setspace}
\setstretch{1.08}
\pagestyle{plain}
\sloppy\sloppypar\raggedbottom\frenchspacing\thispagestyle{empty}

% problem set header
\newcommand{\psname}{Problem Set}
\newcommand{\startps}[1]{\section*{Computational Physics --- \psname~{#1}}}

% problem formatting
\usepackage{amsthm}
\theoremstyle{definition}
\newcommand{\problemname}{Problem}
\newtheorem{problem}{\problemname}
\newcommand{\probref}[1]{\problemname~\ref{#1}}
\theoremstyle{remark}
\newcommand{\notename}{Note}
\newtheorem{note}{\notename}[problem]
\newcommand{\noteref}[1]{\textit{\notename}~\ref{#1}}
\newcounter{subpart}[problem]
\newcommand{\subpart}{\par\addvspace{1ex}\refstepcounter{subpart}\textsl{(\alph{subpart})}~}

% words
\newcommand{\foreign}[1]{\textsl{#1}}
\newcommand{\vs}{\foreign{vs}}

% code
\newcommand{\code}[1]{\texttt{\detokenize{#1}}}
\usepackage{xcolor}
% Define custom colors for code highlighting
\definecolor{codegreen}{rgb}{0,0.6,0}
\definecolor{codegray}{rgb}{0.5,0.5,0.5}
\definecolor{codepurple}{rgb}{0.58,0,0.82}
\definecolor{backcolour}{rgb}{0.99,0.99,0.99}
\usepackage{listings}
\lstdefinestyle{pythonstyle}{
    backgroundcolor=\color{backcolour},   
    commentstyle=\color{codegreen},
    keywordstyle=\color{magenta},
    numberstyle=\tiny\color{codegray},
    stringstyle=\color{codepurple},
    tabsize=4
}
\lstset{
    basicstyle=\small\ttfamily,
    columns=fixed,
    breaklines=false,
    keepspaces=true,
    showspaces=false,
    showstringspaces=false,
    showtabs=false,
    showlines=*,
    style=pythonstyle,
    language=Python,
    literate={~}{{\url{~}}}1
}

% math
\usepackage{amsmath}
\newcommand{\dd}{\mathrm{d}}
\newcommand{\e}{\mathrm{e}}

% primary units
\newcommand{\rad}{\mathrm{rad}}
\newcommand{\kg}{\mathrm{kg}}
\newcommand{\m}{\mathrm{m}}
\newcommand{\s}{\mathrm{s}}

% secondary units
\renewcommand{\deg}{\mathrm{deg}}
\newcommand{\km}{\mathrm{km}}
\newcommand{\cm}{\mathrm{cm}}
\newcommand{\mm}{\mathrm{mm}}
\newcommand{\ft}{\mathrm{ft}}
\newcommand{\mi}{\mathrm{mi}}
\newcommand{\AU}{\mathrm{AU}}
\newcommand{\ns}{\mathrm{ns}}
\newcommand{\h}{\mathrm{h}}
\newcommand{\yr}{\mathrm{yr}}
\newcommand{\N}{\mathrm{N}}
\newcommand{\J}{\mathrm{J}}
\newcommand{\eV}{\mathrm{eV}}
\newcommand{\W}{\mathrm{W}}
\newcommand{\Pa}{\mathrm{Pa}}

% derived units
\newcommand{\mps}{\m\,\s^{-1}}
\newcommand{\mph}{\mi\,\h^{-1}}
\newcommand{\mpss}{\m\,\s^{-2}}

% random units
\newcommand{\au}{\mathrm{a.u.}}


\begin{document}

\startps{3}
Due Friday 2026 September 25 at 18:00 on Brightspace.
Be sure to explicitly give credit where it is due, on all problems.
Also credit buddies, web pages, and LLMs (as per course policy).
Failure to give explicit credits will result in points deducted.

\begin{problem}
  Reconsider the problem of integrating under the semicircle that you did in \psname~2, \problemname~4.
  Choose a sensible value for $N$ and instead of using \code{np.sum()}, do the sum four different ways, listed below.
  Check that all three methods give you the same answer as you got with \code{np.sum()}.
  Which is the most natural of these for your problem?

  For this problem, hand in your code, not just your answers.
  \subpart Construct the sum with a Python \code{for} loop.
  \subpart Construct the sum with a Python \code{while} loop.
  \subpart Construct the sum with a Python \code{map} and \code{reduce}.
\end{problem}

\begin{note}
  The point of this problem is to exercise Python controls, not to do the integral!
\end{note}

\begin{note}
  If you want an extra challenge (not for credit!), construct the sum with \emph{recursion}.
  A more natural use of recursion is the bisection code you wrote in \psname~2.
  The bisection method has a recursive structure; if you want to understand recursion that's another good side quest.
\end{note}

\begin{problem}
  \subpart Newman, \textit{Computational Physics, 2ed}, Exercise 5.2.
  \subpart Now do the same integral again, but using the \code{quadrature} method in the python \code{scipy} package.
  Use the same (or similar) numbers of points or function evaluations.
  \subpart Now make a graph that compares the Simpson's and quadrature integrations as a function of the number of slices (or samples or function evaluations, whichever is easiest).
  What do you conclude?
\end{problem}

\begin{problem}
  Make a diagram or an infographic that shows a smooth, curvy function $y(x)$ as a curve,
  a narrow integration segment $\Delta x$, some relevant function evaluations as dots,
  and the relevant (and different) areas that one is summing in three different integration methods.
  Think of this as an ``infographic'' that explains the difference between the three methods.
  Hogg will say more in Lecture about what he expects.
  For the three methods, one suggestion would be midpoint or Euler's method for the simplest,
  trapezoidal rule for the medium, and Simpson's or second-order for the more sophisticated method.
\end{problem}

\begin{note}
  This is supposed to be something that you might add to a textbook to explain the difference between the three methods.
  Therefore: Label what is important to label!
  But don't include unnecessary axes or numbers or tick marks that aren't used in the explanation.
  You might have learn to control \code{matplotlib} for this \problemname.
\end{note}

\begin{problem}
  Newman, \textit{Computational Physics, 2ed}, Exercise 5.6.
\end{problem}

\end{document}
